% from ega4-8-fr.tex
\oldpage[IV-3]{14}
\item[\emph{c)}] Il existe un $\lambda$ tel que $F_\lambda=S_\lambda$ (et
alors $F_\mu=S_\mu$ pour $\mu\geq\lambda$).
\end{enumerate}
\end{corollary}

Cela se déduit aussitôt de (8.3.3) par passage aux complémentaires.

\begin{corollary}[8.3.5]
\label{IV.8.3.5-fr}
Pour tout $\lambda$, soient $E_\lambda$, $F_\lambda$ deux parties
constructibles de $S_\lambda$ telles que, pour $\mu\geq\lambda$, on ait
$E_\mu\subset u_{\lambda\mu}^{-1}(E_\lambda)$ et
$F_\mu\supset u_{\lambda\mu}^{-1}(F_\lambda)$. Supposons qu'il existe un
indice $\alpha$ tel que $S_\alpha$ soit quasi-compact. Alors, pour que
$E\subset F$ (resp. $E=F$), il faut et il suffit qu'il existe $\lambda$ tel
que $E_\lambda\subset F_\lambda$ (resp. $E_\lambda=F_\lambda$), auquel cas
on a aussi $E_\mu\subset F_\mu$ (resp. $E_\mu=F_\mu$) pour
$\mu\geq\lambda$.
\end{corollary}

C'est un cas particulier de (8.3.2).

En particulier~:

\begin{corollary}[8.3.6]
\label{IV.8.3.6-fr}
Supposons qu'il existe un $\alpha$ tel que $S_\alpha$ soit quasi-compact.
Pour que $S=\emptyset$, il faut et il suffit qu'il existe un $\lambda$ tel
que $S_\lambda=\emptyset$.
\end{corollary}

\begin{corollary}[8.3.7]
\label{IV.8.3.7-fr}
On a, pour tout $\lambda$,
\begin{equation}
u_\lambda(S)=\bigcap_{\mu\geq\lambda}u_{\lambda\mu}(S_\mu).
\tag{8.3.7.1}\label{IV.8.3.7.1-fr}
\end{equation}
\end{corollary}

Il est clair que le premier membre de (8.3.7.1) est contenu dans le second.
Soit $s$ un point de $S_\lambda$ et posons
$X_\lambda=\operatorname{Spec}(k(s))$~; considérons le système projectif
$(X_\mu,v_{\mu\nu})$ où
$X_\mu=X_\lambda\times_{S_\lambda}S_\mu$ et
$v_{\mu\nu}=1\times u_{\mu\nu}$ pour
$\lambda\leq\mu\leq\nu$~; sa limite projective est
$X=X_\lambda\times_{S_\lambda}S$ et $v_\lambda=1\times u_\lambda$ est
l'application canonique $X\to X_\lambda$ (8.2.5). Si
$s\in\bigcap_{\mu\geq\lambda}u_{\lambda\mu}(S_\mu)$, cela entraîne que
$X_\mu\neq\emptyset$ pour tout $\mu\geq\lambda$ (I, 3.4.8)~; il résulte
alors de (8.3.6) que $X\neq\emptyset$, donc $s\in u_\lambda(S)$ par
(I, 3.4.8).

\begin{proposition}[8.3.8]
\label{IV.8.3.8-fr}
\begin{enumerate}
\item[(i)] Pour que le morphisme $u_\lambda:S\to S_\lambda$ soit dominant
(resp. surjectif), il faut et il suffit que pour $\mu\geq\lambda$ le
morphisme $u_{\lambda\mu}:S_\mu\to S_\lambda$ soit dominant (resp.
surjectif).
\item[(ii)] Si, pour un indice $\lambda$, les morphismes
$u_{\lambda\mu}:S_\mu\to S_\lambda$ sont plats (resp. fidèlement plats)
pour tout $\mu\geq\lambda$, alors le morphisme
$u_\lambda:S\to S_\lambda$ est plat (resp. fidèlement plat).
\item[(iii)] Supposons que les morphismes $u_\mu:S\to S_\mu$ soient
surjectifs pour $\mu\geq\lambda$. Pour que $u_\lambda$ soit un morphisme
ouvert (resp. universellement ouvert), il faut et il suffit que, pour tout
$\mu\geq\lambda$, $u_{\lambda\mu}$ soit un morphisme ouvert (resp.
universellement ouvert).
\end{enumerate}
\end{proposition}

\noindent
(i) Comme $u_\lambda(S)\subset u_{\lambda\mu}(S_\mu)$ pour
$\mu\geq\lambda$, la nécessité des conditions est triviale, et il résulte
aussitôt de (8.3.7.1) que si les $u_{\lambda\mu}$ sont surjectifs, il en est
de même de $u_\lambda$. Supposons maintenant les $u_{\lambda\mu}$ dominants
pour $\mu\geq\lambda$, et considérons dans $S_\lambda$ un ouvert
quasi-compact non vide $U_\lambda$~; alors les
$U_\mu=u_{\lambda\mu}^{-1}(U_\lambda)$ pour $\mu\geq\lambda$ forment un
système projectif dont la limite projective est
$U=u_\lambda^{-1}(U_\lambda)$ (8.2.5). Par hypothèse les $U_\mu$ sont tous
non vides, donc il en est de même de $U$ par (8.3.6), ce qui prouve que
$u_\lambda$ est dominant.

\medskip
\noindent
(ii) En vertu de (i), il suffit de considérer le cas où les
$u_{\lambda\mu}$ sont plats~; on peut alors se borner au cas où $S_\lambda$
est affine, donc aussi les $S_\mu$ pour $\mu\geq\lambda$ et $S$, et
l'assertion résulte alors de (2.1.2) et $(0_{\mathrm I},6.2.3)$.

\medskip
\noindent
(iii) En vertu de (8.2.5) et de (I, 3.5.2), il suffit de traiter le cas des
morphismes ouverts. Comme $u_\lambda=u_{\lambda\mu}\circ u_\mu$ et que
$u_\mu$ est surjectif, on sait que si $u_\lambda$ est ouvert il en est de
même de $u_{\lambda\mu}$ pour $\mu\geq\lambda$ (Bourbaki,
\emph{Top. gén.}, chap. $\mathrm I^{\mathrm{er}}$, $3^{\mathrm e}$ éd.,
\S~5, n\textsuperscript{o} 1, prop. 1).

